\documentclass[protokoll]{dvd}

\KOMAoptions{
	headwidth = 18cm,
	footwidth = 18cm,
}

\setlength{\parskip}{0.2em}
% --------------------------------------------------------------------- %
% ---------------------------- Definitions ---------------------------- %
% --------------------------------------------------------------------- %

% Styrelsen
\def\vak{\emph{vakant}}

\def\ordf{Daniel Cole}
\def\vice{\vak}
\def\samo{\vak}
\def\kassa{Aowis Tabbaa}
\def\sekr{\vak}

% mötesinfo

\def\nr{2}
\def\yr{2026}
\def\ymd{2026-09-14} % mötets datum
\def\starttime{17:17}
\def\loc{EA}

% formalia
\def\kall{2026-08-31}

% --------------------------------------------------------------------- %

\begin{document}

\title{Divisionsstämmans möte \nr}
\subtitle{\yr}
\author{Divisionsstämman}
\date{\ymd}

\textbf{Datum:} \csname @date\endcsname\\
\textbf{Tid:} \starttime\\
\textbf{Plats:} \loc\\
\textbf{Styrelsemedlemmar:}
\begin{närvarande_förtroendevalda}
    \förtroendevald{Ordförande}{\ordf}{Ja}
    \förtroendevald{Vice ordförande}{\vice}{---}
    \förtroendevald{SAMO}{\samo}{---}
    \förtroendevald{Kassör}{\kassa}{Ja}
    \förtroendevald{Sekreterare}{\sekr}{---}
\end{närvarande_förtroendevalda}

\textbf{Närvarande övriga medlemmar:}
\begin{närvarande_medlemmar}
    \medlem{Alexander Appelin}
    \medlem{Sarah Manktelow}
    \medlem{Ola Forss}
\end{närvarande_medlemmar}
%\newpage

\section{Öppnande av möte}
Mötet beräknas öppnas av \ordf\ \starttime

Mötet öppnades \starttime

\newpage
\section{Formalia}

\subsection{Fastställande av röstlängd}
En röstlängd är en lista på personer som har rösträtt. Under våra stämmomöten äger endast divisionsmedlemmar som närvarar vid mötet rösträtt. Därmed krävs det att för att stämmomötet ska kunna genomföras behöver vi en lista med namn på röstberättigade medlemmar som närvarar vid stämmosammanträdet.

Medlemmar kan under mötets gång, det vill säga efter denna punkt, läggas till eller tas bort ur röstlängden.
Det ska då framgå i röstlängden vid vilken punkt i mötesschemat medlemmen lämnar eller ankommer till
stämmomötet.

\subsubsection*{Förslag till beslut}
\begin{attsatser}
    \item stämman fastställer den nuvarande röstlängden
\end{attsatser}

\subsubsection*{Beslut}
\begin{attsatser}
    \item attsats 1 bifalles.
\end{attsatser}


\subsection{Divisionsstämmans beslutbarhet}

6 kap. i stadgan definierar regler för Divisionstämman.

\kall \ 
kallade styrelsen till divisionsstämma genom att skriva i föreningens discordserver \emph{MonadenDV}, samt annonserades detta på \emph{dvet}.se \kall

\subsubsection*{Förslag till beslut}

\begin{attsatser}
    \item divisionsstämman har uppnått kraven i stadgan för att få hålla möte, och är därmed beslutbar.
\end{attsatser}

\subsubsection*{Beslut}
\begin{attsatser}
    \item attsats 1 bifalles.
\end{attsatser}


\subsection{Fastställande av mötesschema}

För att divisionsstämman ska kunna fatta ett beslut eller protokollföra en diskussion behöver punkten i mötesschemat där stämman ska fatta beslut vara inlagd eller föras in i mötesschemat senast vid den här punkten.

\subsubsection*{Förslag till beslut}

\begin{attsatser}
    \item mötesschemat fastställs utan några ändringar.
\end{attsatser}

\subsubsection*{Beslut}
\begin{attsatser}
    \item attsats 1 bifalles.
\end{attsatser}

\subsection{Val av mötesordförande}

Mötesordförande har till uppgift att leda Divisionsstämmans sammankomst.
Hen ansvarar för att mötesformalia sköts.



\subsubsection*{Förslag till beslut}
\begin{attsatser}
    \item \ordf\ väljs till mötesordförande.
\end{attsatser}

\subsubsection*{Beslut}
\begin{attsatser}
    \item attsats 1 bifalles.
\end{attsatser}

\subsection{Val av vice mötesordförande}

Vice mötesordförande hjälper mötesordförande med att hålla talarlistan, och att alla får komma till tals.

\subsubsection*{Förslag till beslut}
\begin{attsatser}
    \item \kassa\ väljs till vice mötesordförande
\end{attsatser}

\subsubsection*{Beslut}
\begin{attsatser}
    \item attsats 1 bifalles.
\end{attsatser}

\subsection{Val av mötessekreterare}

Mötessekreteraren har till uppgift att anteckna diskussioner, beslut, och eventuella reservationer under mötet.

\subsubsection*{Fri nominering}


\subsubsection*{Förslag till beslut}
\begin{attsatser}
    \item \ordf\ väljs till mötessekreterare
\end{attsatser}

\subsubsection*{Beslut}
\begin{attsatser}
    \item attsats 1 bifalles.
\end{attsatser}

\subsection{Val av protokolljusterare}

Protokolljusterare har till uppgift att kontrollera att protokollet i slutändan reflekterar de faktiska besluten och diskussionerna som fördes under sammanträdet; samt agera rösträknare vid slutna omröstningar.
Utöver protokolljusterarna så ska mötesordförande och mötessekreteraren signera protokollet.
Vid Divisionsstämmans sammanträden ska det vara två justerare.\\
Mötesordförande och mötessekreteraren kan inte vara justerare.
\subsubsection*{Fri nominering}
Ola Forss och Sarah Manktelow nominerar sig själva. 
\subsubsection*{Förslag till beslut}
\begin{attsatser}
    \item Ola Forss och Sarah  Manktelow väljs till protokolljusterare för denna stämma.
\end{attsatser}

\subsubsection*{Beslut}
\begin{attsatser}
    \item attsats 1 bifalles.
\end{attsatser}

\newpage

\section{Rapporter}

\subsection{Styrelsen}

\subsubsection*{Ordförande}
Sedan senaste stämman har jag skickat in föreningsansökan till Göta samt medverkat på flertalet av sektionens möten. Jag deltog på konstituerande fullmäktigemötet för verksamhetsåret 2026/27 där jag representerade Knapptryckarkompaniet, mötet drog över och vi behövde fortsätta en andra dag och totalt tog det strax över 25 timmar. Under valborgsveckan deltog Divisionen på NaT-seks O'Learys dag och höll en dart-turnering. 
Vi har också skickat in ansökan att ändra fullmäktige gentemot banken, pro-tip för framtiden: gör detta ASAP och börja förbereda innan överlämningen.

Vi blev efterfrågade kompletteringar på båda ansökningar att det behöver vara tre ledamöter i styrelsen. 

Göta vill också att vi ska 
\begin{quote}
     lägga till i stadgarna att det krävs ett ordinarie medlemskap i Göta studentkår för att inneha styrelseuppdrag i föreningen
\end{quote}
samt sammanfatta
\begin{quote}
    vilka rättigheter en medlem har alternativt vad som definierar en medlem.
\end{quote}
Jag håller på att formulera förslag på stadgeändringar för att införa detta. Samt ett förslag att införa minimikrav på tre styrelsemedlemmar. 
Stadgeändringarna behöver vara införda innan nästa föreningsansökan så de är inte lika akut som styrelsemedlemsantalet.


Göta studentkår har nu också kallat till fullmäktigemöte 1 för 2026/-27 den 30:e september. Under sommaren fick också Knapptryckarkompaniet en till ledamot så vi har då två representatner i FUM! 


Under insparken har jag deltagit på alla Götas introdagar samt föreningstorget för att representera Divisionen. 

-- \ordf, Styrelsen 2026

\subsubsection*{Kassör}
Det var en oproduktiv år för mig som kassör eftersom jag har fortfarande ingen tillgång till kontot. Men det var något gjort i alla fall eftersom jag rakade glömma att skicka fakturan för insparksbudgeten men det var också en gång som jag hörde om att skicka en faktura och fick ingen meddelande från Matti gällande insparkspengar.

\subsection{Verksamhetsrapporter av kommittéer}

\subsubsection*{DVRK}
\begin{quote}
    Insparken har gått bra, nu är det bara finalsittningen kvar där vi har beställt catering till maten. Typ alla förberedelser klara och vi ska fixa det sista som är kvar närmare dagen av sittningen. 
    I helhet har allt gått bra, det blev lite flummigt ibland men det blev löst med god kommunikation. 
\end{quote}

\subsubsection*{ConCats}
Vi hade en del framgångsrika arr denna inspark nyligen, de var flertals mer reccar som faktiskt deltog och hade kul som var oförväntat jämforande med förra årets inspark. Förutom det hade jag också pratat med Dlude att ha en gemensam event antingen i basen eller monaden som skulle hända efter insparken. Jag hade också hittat en passande bild för amazing race för förra årets inspark och behöver också ringa vaktmästaren för att fixa toas handfat i monaden.

\subsubsection*{Femme++}
Vi höll två fina arr under insparken Lunchen blev riktigt lyckad med, onAir var lyckad men lite synd att det inte kom dit så många. 

-- rubix

\subsection{Studienämnden}
Jag har varit för upptagen med annat för att göra något konkret, jag har en idé för mitt-i-kursen pluggstuga+kursutvärdering så kanske kommer något sådant om ett par veckor. 

-- Cassilda

\subsubsection*{Mega6}
Vi hade ett krök, gick bra. Vi hade en sittning, gick jättebra. Bara det.

\subsubsection*{DVOps}
Dv\textunderscore ops har sedan förra stämman gjort lite för en gångs skull. Vi programmerade äntligen klart discordbotten “Monadic Mission” har släppt två problem och vi kommer att släppa mer problem efter insparken. Under insparken höll jag eventet “Scratch Game Jam”, där många riktigt bra spel skapades och det var kul. Vill gärna ha mer hackathon events i framtiden. 

-Wildcard


\subsubsection*{Mega7}
Ingen rapport.

\subsubsection*{Klubbsport DV}
Ingen rapport.


\newpage

\section{Beslutspunkter}

Enligt Stadgan måste ändringar av Stadgan röstas igenom på två av Divisionsstämmans varandra följande möten.
Om en beslutspunkt innehåller ``första läsningen'' innebär det att det är första gången beslutet tas upp för omröstning.
Om en beslutspunkt innehåller ``andra läsningen'' innebär det att beslutspunkten har röstats igenom förra stämmomötet, och beslutet behöver bekräftas för att gå igenom.


\textit{Inga beslutspunkter eller motioner har inkommit innan stämman.}

%\newpage

\section{Inval}
Inval av personer till förtroendeposterna som väljs in av Divisionsstämman. Dessa väljs in för en ordinarie mandatperiod, vilket sträcker sig från 1 januari till 31 december.

\subsection{Styrelsen}
\subsubsection*{Inkomna nomineringar inför mötet}
Angelika har varit en av föreningens mest engagerade medlemmar under de senaste åren. Med flera år av erfarenhet inom Mega6 tycker vi hon är mycket lämpad till posten som SAMO för den resterande mandatperioden.

Vi i Styrelsen '26 vill därför nominera Angelika Hagberg till SAMO för Styrelsen '26. 

\subsubsection*{Förslag till beslut}
\begin{attsatser}
    \item Angelika Hagberg väljs till SAMO för Styrelsen '26
\end{attsatser}
%\subsubsection*{Utjustering}
%\textit{Kandidaten} justeras ut.
\subsubsection*{Beslut} 
\begin{attsatser}
    \item attsats 1 bifalles.
\end{attsatser}
%\subsubsection*{Injustering}
%\textit{Kandidaten} justeras in.

\subsection{DVArm}
\subsubsection*{Inkomna nomineringar inför mötet}
\textit{Styrelsen har ej fått in några nominering innan stämman.}
\subsubsection*{Fri nominering}
\subsubsection*{Förslag till beslut}
\begin{attsatser}
    \item punkten bordsläggs
\end{attsatser}
\subsubsection*{Beslut}
\begin{attsatser}
    \item attsats 1 bifalles.
\end{attsatser}

\subsection{Revisor}
5 kap, 3 § i stadgan definierar följande krav
\begin{quote}
    Divisionens revisorer kan inte inneha något annat förtroendeuppdrag i divisionen under dennes mandatperiod som revisor.\\
    
    Revisorerna måste även vara myndiga och inte vara försatta i konkurs. De behöver även ha tillräckliga kunskaper i redovisning och ekonomiska förhållanden för att kunna utföra sina uppdrag.
\end{quote}
\subsubsection*{Inkomna nomineringar inför mötet}
\textit{Styrelsen har ej fått in några nominering innan stämman.}
\subsubsection*{Fri nominering}
\subsubsection*{Förslag till beslut}
\begin{attsatser}
    \item punkten bordsläggs
\end{attsatser}
\subsubsection*{Beslut}
\begin{attsatser}
    \item attsats 1 bifalles.
\end{attsatser}

%\newpage
\section{Diskussionspunkter}
Stämman är inte bara en chans för oss i divisionen att rösta om saker, utan den ger oss även en chans att diskutera olika ämnen, som kanske nödvändigtvis inte behövs röstas om.

Styrelsen har ej kommit med några diskussionspunkter, men vi lämnar golvet öppet för närvarande medlemmar att ta upp det dom har på hjärtat.\\

\textit{Inga diskussionspunkter har inkommit innan mötet, och styrelsen har ej kommit några punkter själva.}
%\subsection*{Eventuella diskusionspunkter}

\section{Avslutande av möte}

Mötet avslutades klockan 17:59

\stämmosignaturer

\end{document}
